\documentclass[11pt,a4paper]{article}
\usepackage[T1]{fontenc}
\usepackage[utf8]{inputenc}
\usepackage{lmodern}
\usepackage[margin=28mm]{geometry}
\usepackage{amsmath,amssymb,amsthm,mathtools,mathrsfs}
\usepackage{microtype}
\usepackage{booktabs,longtable,tabularx,array}
\usepackage{enumitem}
\usepackage{graphicx,xcolor,tikz}
\usetikzlibrary{arrows.meta,positioning,calc,decorations.pathreplacing,patterns}
\usepackage[hyphens]{url}
\usepackage[colorlinks=true,linkcolor=blue!45!black,citecolor=blue!45!black,
 urlcolor=blue!45!black,hypertexnames=false]{hyperref}
\usepackage{bookmark}
\allowdisplaybreaks[2]
\setlength{\emergencystretch}{3em}
\setlength{\parskip}{3pt}
\setcounter{tocdepth}{1}
\setlist{itemsep=2pt,topsep=4pt,parsep=0pt}
\setlist[enumerate,1]{label=\textup{(\roman*)}}

% ---------------------------------------------------------------- theorems
\newtheorem{theorem}{Theorem}[section]
\newtheorem{lemma}[theorem]{Lemma}
\newtheorem{proposition}[theorem]{Proposition}
\newtheorem{corollary}[theorem]{Corollary}
\newtheorem{maintheorem}{Theorem}
\renewcommand{\themaintheorem}{\Alph{maintheorem}}
\theoremstyle{definition}
\newtheorem{definition}[theorem]{Definition}
\newtheorem{example}[theorem]{Example}
\newtheorem{setup}[theorem]{Standing assumptions}
\theoremstyle{remark}
\newtheorem{remark}[theorem]{Remark}

% A visible but light marker for the steps of a long proof.
\newcommand{\step}[1]{\par\smallskip\noindent\textit{Step #1.}\ }

% ---------------------------------------------------------------- general
\newcommand{\R}{\mathbb R}
\newcommand{\T}{\mathbb T}
\newcommand{\Z}{\mathbb Z}
\newcommand{\av}[1]{\langle #1\rangle}
\newcommand{\dd}{\,d}
\newcommand{\id}{\mathrm{id}}
\newcommand{\supp}{\operatorname{supp}}
\newcommand{\dist}{\operatorname{dist}}
\newcommand{\diag}{\operatorname{diag}}
\newcommand{\Ai}{\operatorname{Ai}}
\newcommand{\Bi}{\operatorname{Bi}}

% ---------------------------------------------------------------- notation
% Symbols whose letters would otherwise collide are given macros, so
% that each has one meaning throughout the paper.
\newcommand{\jn}{\mathsf j}             % band edges  j_n = C_* q_n A_n^{-1/2}
\newcommand{\ea}{\mathsf a}             % coefficient exponent of sources
\newcommand{\eb}{\mathsf b}             % coefficient exponent of outputs
\newcommand{\el}{\mathsf l}             % intermediate loss exponent (Section 6)
\newcommand{\bd}{\mathsf d}             % boundary-defining function y(1-y)
\newcommand{\Dr}{\mathsf D}             % unit-scale radial derivative
\newcommand{\Qint}{\mathsf Q}           % integral functional on sources
\newcommand{\Odd}{\mathsf O}            % odd projection on the circle
\newcommand{\Even}{\mathsf E}           % even projection on the circle
\newcommand{\refl}{\mathsf s}           % reflection theta -> -theta
\newcommand{\Ev}{\mathsf{Ev}}           % even substitution at a root
\newcommand{\Od}{\mathsf{Od}}           % odd divided substitution at a root
\newcommand{\mA}{\mathsf A}             % matching matrices of Section 6
\newcommand{\mB}{\mathsf B}
\newcommand{\mK}{\mathsf K}
\newcommand{\mS}{\mathsf S}
\newcommand{\bx}{\mathbf x}             % the state (u,nu,kappa,h_L)
\newcommand{\bW}{\mathbf W}             % augmented canonical unknown

\newcolumntype{Y}{>{\raggedright\arraybackslash}X}
\hypersetup{pdftitle={Smooth isometric bending of a closed sphere},
 pdfauthor={Boyang Hu}}
